\section{The exact two-leg specialization}
Original polynomial and marked points: The Clankers \cite{ES488}.


This note constructs two explicit target paths for the original
ES--Fable polynomial. Their concatenated specialization has
eight, then seven, then three finite points in the original
source chart. It is not the single arithmetic trajectory in
MO12 of the preceding sections: both paths and all their signs are
specified here. In particular its second-leg escape exponent
differs from that transverse arithmetic trajectory.

\section{The original map and its exact fibre coordinates}
Use the original source coordinates $(a,y,z,w)$. Its complete
factor-coordinate formulas and polynomial map are
\begin{align}
 b&=i+ay,&c&=-i+2ay+a^2z,\nonumber\\
 d&=-iy-a(iz+2y^2)-a^2yz,&e&=2z-7iy^2+aw,\nonumber\\
 f&=iw+3iy^3-4yz+a(6iy^2z+wy+4y^4)+2a^2y^3z,&&\nonumber\\
 P(a,y,z,w)&=(ac,ae+bd,af+be,bf).
 &&\tag{VA1}
\end{align}
Here $L=aU+bV$, $C=cU^3+dU^2V+eUV^2+fV^3$ satisfy
$ad+bc=1$ and $C(-b,a)=1$, directly by substitution. Thus
their product at target $u=(A,B,C_1,D)$ is
$H_u=AU^4+U^3V+BU^2V^2+C_1UV^3+DV^4$.
The symbol $C_1$ is the third target coefficient, distinct
from the cubic factor denoted $C$. No source coordinate or
polynomial coefficient has been changed in VA1.

For a simple finite root $r$ of
$h_u(x)=Ax^4+x^3+Bx^2+C_1x+D$, the two inverse states are
\begin{align}
 a^2&=1/h_u'(r),&b&=-ar,&c&=A/a,&d&=(1+Ar)/a,\nonumber\\
 e&=(B+r+Ar^2)/a,& f&=(C_1+Br+r^2+Ar^3)/a,&&\nonumber\\
 y&=-r-i/a,&z&=A/a^3+2r/a+3i/a^2,&&\nonumber\\
 w&=7ir^2/a+(B-17r+Ar^2)/a^2-13i/a^3-2A/a^4.
 &&\tag{VA2}
\end{align}
Both nonzero square roots for $a$ are retained. To prove
completeness, divide $H_u$ by $a(U-rV)$ and compare successive
coefficients to obtain the six factor entries in VA2.
The resultant is $C(-b,a)=a^2h_u'(r)$, which must be one;
thus a multiple root gives no state. The original chart inverse
$y=(b-i)/a$, $z=(c+i-2ay)/a^2$,
$w=(e-2z+7iy^2)/a$ gives the remaining entries of VA2.
All operations are reversible for $a\ne0$.
At $A=0$ the root at infinity is simple, because the original
cubic coefficient is one. Its positive-chart state is exactly
\[
 q_\infty(A=0,B,C_1,D)
 =\left(0,B,\frac{i(7B^2-C_1)}2,11B^3-2BC_1-D\right).
\]
Indeed VA1 at $a=0$ gives
$P=(0,y,2iz+7y^2,-4iyz-w-3y^3)$, which solves to this formula.
The opposite factor sign has $b=-i$ and is outside the original
chart, since $a=0$ forces $b=i$ in VA1. Therefore every
simple finite root contributes two source states and infinity
contributes one in this chart.

\section{First leg: eight states specialize to the original seven}
For $0<\epsilon<1/2$ put $m_\epsilon=1-\epsilon^2/2$ and
$d_\epsilon=\sqrt{4-9\epsilon^2/4}>0$. The first target path,
its factorization and its four distinct marked roots are
\begin{align}
 u^{(1)}_\epsilon
   &=(-i\epsilon+i\epsilon^3/2,3i\epsilon/2,-1,0)
                \longrightarrow h_0=(0,0,-1,0),\nonumber\\
 h^{(1)}_\epsilon(x)
   &=(\epsilon x+i)x
             \left(-im_\epsilon x^2+\frac\epsilon2 x+i\right),\nonumber\\
 r_{\mathrm M}&=-i/\epsilon,&r_{\mathrm N}&=0,&
 r_{\mathrm T}&=\frac{d_\epsilon-i\epsilon/2}{2m_\epsilon},&
 r_{\mathrm E}&=\frac{-d_\epsilon-i\epsilon/2}{2m_\epsilon},\nonumber\\
 \operatorname{Disc} H_{u^{(1)}_\epsilon}&=4-9\epsilon^2/4>0.
 &&\tag{VA3}
\end{align}
The factorization follows by multiplying its displayed factors.
The quadratic formula gives the last two roots. The full
quartic discriminant is obtained from the resultant of this
polynomial and its derivative; substitution and expansion
gives exactly the last line of VA3. Its leading coefficient
$-i\epsilon m_\epsilon$ is nonzero on the stated interval.
Consequently all four roots are finite and distinct and VA2
gives exactly eight source states, with no root collision
anywhere before the endpoint. Their full inverse coordinates
are VA2 evaluated at the four explicit roots VA3 and both
values of $a^2=1/(h^{(1)}_\epsilon)'(r)$.

For the M root these two original-coordinate states simplify
completely to
\[
 q_{\mathrm M,+}(\epsilon)=(\epsilon,0,i/2,0),\qquad
 q_{\mathrm M,-}(\epsilon)
 =(-\epsilon,2i/\epsilon,6i/\epsilon^2-i/2,40i/\epsilon^3).
\]
Substitution in VA1 proves both have image VA3. Equivalently
their selected derivatives are $\epsilon^{-2}$ and their
factor scalings are $a=\pm\epsilon$. Their factor states
are negatives of each other. The positive one tends to
$(0,i,-i,0,i,0)$; the negative one tends to
$(0,-i,i,0,-i,0)$, a finite point in the factor completion
excluded from the original chart. Only its chart coordinates
escape. For the three remaining roots, VA3 gives limits
$0,1,-1$ and nonzero derivative limits $-1,2,2$.
Every expression in VA2 then has a finite limit on each
retained square-root branch. Thus exactly one of the eight
source states escapes, and seven remain finite at $h_0$.

\section{The complete signed marking at the junction}
The four original displayed collision points are
\[
\begin{aligned}
 q_{\mathrm M}&=(0,0,i/2,0),&q_{\mathrm N}&=(i,-1,-3i,13),\\
 q_{\mathrm T}&=(1/\sqrt2,-1-\sqrt2i,2\sqrt2+6i,-34-19\sqrt2i),\\
 q_{\mathrm E}&=(1/\sqrt2,1-\sqrt2i,-2\sqrt2+6i,34-19\sqrt2i).
\end{aligned}
\]
Define the original linear involution
$J(a,y,z,w)=(-a,-y,z,-w)$. The exact marked factor-root
associations and factor states at $h_0$ are
\begin{equation}
\begin{array}{c|c|c|c}
\text{label}&\text{root}&\text{positive factor state }(a,b,c,d,e,f)
                         &\text{two source-chart states}\\\hline
\mathrm M&\infty&(0,i,-i,0,i,0)&q_{\mathrm M},\ \text{opposite sign omitted}\\
\mathrm N&0&(i,0,0,-i,0,i)&q_{\mathrm N},\ Jq_{\mathrm N}\\
\mathrm T&1&(1/\sqrt2,-1/\sqrt2,0,\sqrt2,\sqrt2,0)
                       &q_{\mathrm T},\ Jq_{\mathrm E}\\
\mathrm E&-1&(1/\sqrt2,1/\sqrt2,0,\sqrt2,-\sqrt2,0)
                       &q_{\mathrm E},\ Jq_{\mathrm T}
\end{array}\tag{VA4}
\end{equation}
The negative factor state in each row is the negative of
every one of its six displayed entries. Multiplying $L$ and
$C$ in every row gives $UV(U-V)(U+V)$, with resultant one.
Substituting the original source states in VA1 recovers these
rows and their negatives. This proves all associations in
VA4. In particular $J$ interchanges the T/E factor-root
marking; $Jq_{\mathrm T}$ is not the T negative-sign state.
The missing M negative factor state is the one displayed
at the end of the preceding section. Thus all eight signed
factor states are described, and exactly seven of them
belong to the original source chart.

\section{Second leg: seven states specialize to three finite states}
Take $0\le s<1$ and put $\tau=1-s\in(0,1]$. The second
target path is
$u^{(2)}(s)=(0,s,s-1,0)=(0,1-\tau,-\tau,0)$.
It starts at $h_0$ when $s=0$ and tends to
$u_*=(0,1,0,0)$ when $s\uparrow1$. Its complete polynomial
and all seven inverse states are
\begin{align}
 h^{(2)}_\tau(x)&=x(x-\tau)(x+1),&
 \operatorname{Disc} H_{u^{(2)}}&=\tau^2(1+\tau)^2>0,\nonumber\\
 q_{\mathrm M}(\tau)&=
 \left(0,1-\tau,\frac{i(7(1-\tau)^2+\tau)}2,
                      11(1-\tau)^3+2\tau(1-\tau)\right),\nonumber\\
 q_{\mathrm N,\sigma}(\tau)&=
 \left(\frac{\sigma i}{\sqrt\tau},-\sigma\sqrt\tau,-3i\tau,
                     -\tau+\tau^2+13\sigma\tau^{3/2}\right),\nonumber\\
 q_{\mathrm T,\sigma}(\tau)&=
 \left(\frac{\sigma}{\sqrt{\tau(1+\tau)}},
       -\tau-i\sigma\sqrt{\tau(1+\tau)},\right.\nonumber\\
 &\hspace{12mm}2\sigma\tau^{3/2}\sqrt{1+\tau}+3i\tau(1+\tau),
 \left.\tau(1+\tau)(1-18\tau)
      -i\sigma\tau^{3/2}\sqrt{1+\tau}(13+6\tau)\right),\nonumber\\
 q_{\mathrm E,\sigma}(\tau)&=
 \left(\frac{\sigma}{\sqrt{1+\tau}},1-i\sigma\sqrt{1+\tau},
       -2\sigma\sqrt{1+\tau}+3i(1+\tau),\right.\nonumber\\
 &\hspace{12mm}\left.(18-\tau)(1+\tau)
               -i\sigma\sqrt{1+\tau}(6+13\tau)\right),
 \qquad\sigma\in\{1,-1\}.\tag{VA5}
\end{align}
All real square roots in VA5 are the positive ones; $\sigma$
retains both actual signs. The finite roots are precisely
$0,\tau,-1$ and their derivatives are $-\tau,\tau(1+\tau),1+\tau$.
They are pairwise distinct for every $0<\tau\le1$, and
their squared pair differences give the discriminant in
VA5. Substitution of these roots and both scalings into
VA2 gives every displayed finite-root formula. The infinity
formula after VA2 gives $q_{\mathrm M}(\tau)$.
Thus VA5 is the complete fibre and has no other root
coalescence before its endpoint.

At $\tau=1$ these formulas agree exactly with VA4:
the N positive and negative states are $q_{\mathrm N}$ and
$Jq_{\mathrm N}$; the T states are $q_{\mathrm T}$ and
$Jq_{\mathrm E}$; the E states are $q_{\mathrm E}$ and
$Jq_{\mathrm T}$. At $\tau=0$ the N and T roots collide.
Their four original source states escape, since their
$a$ coordinates diverge. The three remaining finite limits are
\[
 (0,1,7i/2,11),\qquad
 (\sigma,1-i\sigma,-2\sigma+3i,18-6i\sigma),
 \qquad\sigma\in\{1,-1\}.
\]
The complete fibre rule proved after VA2 gives exactly these
three states at the endpoint: the double root zero gives
none, the simple root $-1$ gives two, and infinity gives
the single retained chart state. The N and T factor
coefficients $a$ themselves diverge, so these four escapes
persist in the two-chart factor completion. The opposite
infinity sign was already outside the original chart all
along this second leg and is not a further escape on it.

\section{Exact fixed original-Gram transfer of the two legs}
Let $F:\C^4\to V$ be any one of the already constructed
fixed injective original receiving maps, with its unchanged
physical Hermitian inner product. Put $G=F^*H_VF>0$.
For example $F$ may be the fixed original conductor map
$\Psi_*$, or the actual injective marked observer $F_B$
on its proved nonzero-corner domain. Its full Gram, including
off-diagonal terms and original mass, is retained below.
Write $e_a=(1,0,0,0)^{\mathsf T}$,
$e_y=(0,1,0,0)^{\mathsf T}$,
$e_z=(0,0,1,0)^{\mathsf T}$ and
$e_w=(0,0,0,1)^{\mathsf T}$. The exact limiting vectors
and norm constants are
\begin{align}
 \epsilon^3q_{\mathrm M,-}(\epsilon)&\longrightarrow40ie_w,&
 \epsilon^3\|Fq_{\mathrm M,-}(\epsilon)\|
                       &\longrightarrow40\sqrt{G_{44}},\nonumber\\
 \sqrt\tau\,q_{\mathrm N,\sigma}(\tau)&\longrightarrow\sigma i e_a,&
 \sqrt\tau\,q_{\mathrm T,\sigma}(\tau)&\longrightarrow\sigma e_a,\nonumber\\
 \sqrt\tau\,\|Fq_{\mathrm N,\sigma}(\tau)\|
    &\longrightarrow\sqrt{G_{11}},&
 \sqrt\tau\,\|Fq_{\mathrm T,\sigma}(\tau)\|
    &\longrightarrow\sqrt{G_{11}},\nonumber\\
 \|F(u^{(2)}(s)-u_*)\|
   &=\tau\sqrt{(e_y+e_z)^*G(e_y+e_z)},&&\nonumber\\
 \|F(u^{(2)}(s)-u_*)\|^{1/2}\|Fq_{l,\sigma}(\tau)\|
   &\longrightarrow
     \bigl((e_y+e_z)^*G(e_y+e_z)\bigr)^{1/4}\sqrt{G_{11}},
 &&l\in\{\mathrm N,\mathrm T\}.
 \tag{VA6}
\end{align}
The vector limits are immediate from the complete coordinates
above. Applying the fixed linear $F$ and continuity of its
original norm proves each limit, with positive constants
because $F$ is injective. The target-distance equality is
exact at every $\tau$, since
$u^{(2)}(s)-u_*=-\tau(e_y+e_z)$.
In particular its coefficient is
$G_{22}+G_{33}+2\Re G_{23}$ under the square root;
the off-diagonal metric term has not been deleted.

These statements prove an exact two-leg specialization
$8\to7\to3$, including its missing factor sign and its
four original root-collision escapes. They do not claim a
single flow trajectory traverses those two paths. The
literal arithmetic orbit through $h_0$ has seven entire
finite trajectories, as proved in MO11 of the preceding sections.
The separate transverse arithmetic orbit MO12--19 has
four quarter-power escapes, one cubic chart escape and
three finite limits at $u_*$. The second path VA5 is
tangent to the discriminant: its discriminant vanishes
quadratically in $\tau$ and its four escape rates are
$\tau^{-1/2}$, exactly as VA6 proves. These are compatible
specializations of the same original polynomial, with
their different paths and rates retained explicitly.

The exact certificate
\texttt{verify\_exact\_marked\_specialization.py} verifies
the two first-leg M states, both discriminants, and all
seven original-coordinate second-leg formulas. For the
latter it uses $t=\sqrt\tau$ and $b=\sqrt{1+\tau}$ and
reduces every polynomial numerator modulo $b^2-t^2-1$
over $\mathbb Q(i)$; no numerical sample replaces an
identity. Its captured output is
\texttt{verify\_exact\_marked\_specialization\_output.txt}.

